% =====================================================================
% III Juegos Florales Agustinos por los Derechos Humanos
% Categoría: Artículo de opinión

\documentclass[10pt,a4paper]{article}

% --- Idioma y codificación -------------------------------------------
\usepackage[spanish,es-noshorthands]{babel}
\usepackage[utf8]{inputenc}
\usepackage[T1]{fontenc}
\usepackage{csquotes}
\usepackage{tabularx}
\usepackage{scrhack}
\usepackage{float}

% --- Página y tipografía ---------------------------------------------
\usepackage[a4paper,margin=2.54cm]{geometry}
\usepackage{newtxtext,newtxmath}
\usepackage{setspace}
\doublespacing
\setlength{\parindent}{1.27cm}
\setlength{\parskip}{0pt}

\usepackage{graphicx}
\usepackage{xcolor}
\usepackage{enumitem}
\usepackage{booktabs}
\usepackage{fancyhdr}

% --- Hipervínculos y bibliografía (hyperref antes de biblatex) -------
\usepackage[hidelinks]{hyperref}
\hypersetup{
  pdftitle={Si no es mi familia, no es mi problema: la indiferencia selectiva del universitario},
  pdfauthor={Camilo Bladimir Alexander Capatinta Almanza}
}
\usepackage[style=apa, backend=biber, sortcites=true]{biblatex}
\DeclareLanguageMapping{spanish}{spanish-apa}
\addbibresource{referencias.bib}

% --- Macros ----------------------------------------------------------
% Marcador de contenido pendiente (se ve en rojo; búscalo con Ctrl+F "[").
\newcommand{\relleno}[1][Contenido pendiente]{\textcolor{red}{\textbf{[#1]}}}

\newcommand{\titulo}{Si no es mi familia, no es mi problema:\\ la indiferencia selectiva del universitario}
\newcommand{\autor}{Camilo Bladimir Alexander Capatinta Almanza}
\newcommand{\autorcorto}{Capatinta Almanza}

% --- Encabezado y pie de página --------------------------------------
\setlength{\headheight}{14.5pt}
\pagestyle{fancy}
\fancyhf{}
\renewcommand{\headrulewidth}{0pt}
\fancyfoot[L]{\small \autorcorto}
\fancyfoot[C]{\small III Juegos Florales Agustinos}
\fancyfoot[R]{\small Página \thepage}

\begin{document}

% =====================================================================
% PORTADA
% =====================================================================
%TC:ignore
\begin{titlepage}
\thispagestyle{empty}
\begin{singlespace}
\centering

\IfFileExists{Escudo_UNSA.png}%
  {\includegraphics[height=3cm]{Escudo_UNSA.png}}%
  {\relleno[Falta Escudo\_UNSA.png en la carpeta]}\\[1em]

{\large\bfseries UNIVERSIDAD NACIONAL DE SAN AGUSTÍN DE AREQUIPA}\\[0.4em]
{\large Defensoría Universitaria}\\[3em]

{\large\bfseries III JUEGOS FLORALES AGUSTINOS\\ POR LOS DERECHOS HUMANOS}\\[0.4em]
{\itshape Entre colores y voces}\\[1.5em]
{\large Categoría: Artículo de opinión}

\vfill

{\LARGE\bfseries \titulo}

\vfill

\begin{tabular}{rl}
  \textbf{Autor:} & \autor \\
  \textbf{Escuela Profesional:} & Ingeniería de Sistemas \\
\end{tabular}

\vfill

Arequipa, Perú\\
2026

\end{singlespace}
\end{titlepage}
\setcounter{page}{1}
%TC:endignore

\begin{center}
{\large\bfseries \titulo}\\[0.5em]
{\normalsize \autor}
\end{center}

\section*{Introducción}
% Objetivo (~110 palabras): gancho, tesis en una frase y anuncio de los tres argumentos.
Defendemos los derechos humanos con una condición que casi nunca decimos en
voz alta: que el golpe nos alcance. Los estudiamos en cursos que se aprueban y
se olvidan, y solo nos indignamos cuando alguien de casa resulta mordido. Como
dice el dicho, la víbora solo ataca al descalzo, y nosotros solo ayudamos al
descalzo que lleva nuestro apellido.
\relleno[Verificar la formulación exacta del dicho antes de usarlo]
\relleno[Tesis en una frase: la indiferencia del universitario no nace de la falta de información, sino de un privilegio que confundimos con normalidad. Anunciar los tres argumentos: psicológico, sociológico-político y filosófico]

\section*{La cercanía como filtro}
% Dimensión psicológica (~150 palabras).
\textcite{darley1968} comprobaron que, mientras más testigos de una emergencia
creían que había, menos probable y más lenta era la ayuda, y propusieron que
cada testigo siente que la responsabilidad se reparte y que alguien más actuará.
Algo parecido ocurre con el sufrimiento anónimo: \textcite{small2003} hallaron
que basta saber quién es la víctima, incluso sin datos personales, para que
aumente la disposición a ayudar.
\relleno[Aplicar al universitario ante una marcha o un reclamo ajeno. Aclarar que son hallazgos de laboratorio usados como analogía y que el efecto de la víctima identificable suele ser pequeño en revisiones posteriores]

\section*{Cursos de relleno y protestas bajo sospecha}
% Dimensión sociológico-política (~180 palabras). Aquí va el dato.
El problema es persistente. En mayo de 2026, la Defensoría del Pueblo
registró 195 conflictos sociales en el país, 153 de ellos activos
\parencite{defensoria2026}. Y quien reclama puede cargar además con un
estigma: la Comisión Interamericana de Derechos Humanos dedicó un apartado de
su informe sobre las protestas en el Perú a la estigmatización
\parencite{cidh2023}.
\relleno[Cuidar: no todo conflicto social es una violación de derechos humanos, y el informe de la CIDH fue discutido en el Perú; escribir \enquote{la CIDH sostiene}. Añadir evidencia propia sobre el \enquote{curso de relleno} (plan de estudios o mini encuesta) y desarrollar el privilegio de no ser afectado]

\section*{Un derecho que solo defiendo cuando me afecta}
% Dimensión filosófica (~150 palabras).
La Enciclopedia del Holocausto del Museo Conmemorativo del Holocausto de
Estados Unidos recuerda a Martin Niemöller, un pastor luterano que simpatizó
con ideas nazis antes de oponerse al régimen y cuya frase sobre el silencio
ante las persecuciones se exhibe en el museo desde 1993
\parencite{niemoeller_ushmm}. Por su parte, \textcite{cortina2017} llama
aporofobia al rechazo al pobre, a quien se supone incapaz de dar algo a cambio.
\relleno[Parafrasear la idea de Niemöller sin copiar el poema (existen varias versiones; cita la que uses). Conectar con quién \enquote{merece} nuestra indignación y cerrar con la idea central: un derecho que solo defiendo cuando me afecta no es un derecho, es un privilegio]

\section*{¿Y si de nada sirve?}
% Objeción y respuesta (~60 palabras).
\relleno[Objeción breve: \enquote{yo solo no cambio nada} o \enquote{no tengo tiempo}. Respuesta en pocas líneas, sin sermonear]

\section*{Conclusión: atreverse}
% Conclusión con reflexión final y propuesta (~120 palabras).
Cambiar de conciencia es lo más básico que se espera de cualquier persona con
conciencia social y de clase, pero hoy ya no alcanza. Hace falta cuestionar,
atreverse y exigir, por cauces pacíficos y legales, frente a lo que pretende
mantenernos cómodos y callados.
\relleno[Concretar \enquote{el sistema}: ¿quién controla al universitario y cómo? Poner primero el principio (la dignidad ajena no depende de que me alcance) y después la advertencia de que lo que se normaliza con uno puede repetirse con otros. Cerrar con uno o dos gestos concretos: preguntar, exigir un curso serio de derechos humanos, organizarse, acompañar con información]

% =====================================================================
% REFERENCIAS
% =====================================================================
%TC:ignore
\clearpage
\printbibliography[title={Referencias}]

% =====================================================================
% ANEXOS
% =====================================================================
\clearpage
\appendix
\renewcommand{\thesection}{Anexo \Alph{section}}

\section{Ficha de datos del participante}
% Datos obligatorios (bases, sección VIII). Se adjuntan junto con el trabajo.
\begin{singlespace}
\renewcommand{\arraystretch}{1.8}
\noindent
\begin{tabularx}{\textwidth}{@{}p{5.2cm}X@{}}
  \toprule
  \textbf{Dato} & \textbf{Información} \\
  \midrule
  Nombres y apellidos completos & \autor \\
  Número de DNI & \relleno[DNI] \\
  Código Universitario (CUI) & \relleno[CUI] \\
  Número de teléfono de contacto & \relleno[Teléfono] \\
  Correo electrónico institucional & \relleno[Correo @unsa.edu.pe] \\
  Escuela Profesional & Ingeniería de Sistemas \\
  Categoría en la que participa & Artículo de opinión \\
  \bottomrule
\end{tabularx}
\end{singlespace}

\bigskip
\noindent\textit{Nota: la omisión de alguno de los datos puede impedir la validación de la participación (bases, sección VIII).}

\section{Declaración de originalidad y autorización}
% Opcional: las bases indican que el envío ya implica aceptación. Bórrala si no la quieres.
\begin{singlespace}
Yo, \autor, estudiante de la Escuela Profesional de Ingeniería de Sistemas de
la Universidad Nacional de San Agustín de Arequipa, declaro que el artículo
\enquote{Si no es mi familia, no es mi problema: la indiferencia selectiva del
universitario} es de mi autoría y creación original, y que no ha sido copiado
ni plagiado ni vulnera derechos de autor de terceras personas.

Declaro conocer y aceptar las bases de los III Juegos Florales Agustinos por
los Derechos Humanos, y autorizo a la Universidad Nacional de San Agustín de
Arequipa, a través de la Defensoría Universitaria, a utilizar, reproducir y
difundir el material presentado con fines exclusivamente institucionales,
educativos, culturales y de difusión, sin fines comerciales. Esta autorización
no implica cesión de mis derechos de autor.

\vspace{3em}
\noindent
\begin{tabular}{@{}p{7cm}p{7cm}@{}}
  Arequipa, \relleno[fecha] & \rule{6cm}{0.4pt} \\
   & \autor \\
\end{tabular}
\end{singlespace}

\section{Matriz de fuentes por dimensión}
% Material de trabajo: ELIMINAR este anexo antes de enviar.
\begin{singlespace}
\small
\renewcommand{\arraystretch}{1.3}
\noindent
\begin{tabularx}{\textwidth}{@{}p{2.3cm}p{3cm}XX@{}}
  \toprule
  \textbf{Dimensión} & \textbf{Fuente} & \textbf{Qué respalda} & \textbf{Matiz a declarar} \\
  \midrule
  Psicológica & \textcite{darley1968} & Más testigos percibidos, menos y más lenta ayuda; difusión de la responsabilidad. & Experimento de laboratorio sobre emergencias; usar como analogía. \\
  Psicológica & \textcite{small2003} & Una identificabilidad mínima aumenta la disposición a ayudar. & Efecto pequeño en revisiones posteriores; no presentarlo como ley. \\
  Sociológico-política & \textcite{defensoria2026} & Persistencia de los conflictos sociales (195 en mayo de 2026; 153 activos). & Mide conflictos, no actitudes; no todo conflicto es una violación de derechos humanos. Confirmar la cifra en el PDF. \\
  Sociológico-política & \textcite{cidh2023} & Apartado sobre estigmatización en el contexto de las protestas. & Informe discutido en el Perú; citarlo como postura de la CIDH. Añadir URL oficial. \\
  Filosófica & \textcite{niemoeller_ushmm} & Origen y contexto de la frase de Niemöller. & Hay varias versiones y se ha cuestionado la del museo; no copiar el poema. \\
  Filosófica & \textcite{cortina2017} & Aporofobia: el rechazo al pobre. & Trata de la pobreza, no de la protesta; conectar con cuidado. \\
  \bottomrule
\end{tabularx}
\end{singlespace}
%TC:endignore

\end{document}